% ============================================================
%  Modular WMC for Neuro-Symbolic Inference — JAIR Submission
%  Revised from AAAI-27 rejection; targeting JAIR 2027
%  Revision addresses: c36U (GO cc degeneracy), KEL1 (WMC
%  solver comparison, complexity), AI Reviewer (related work)
% ============================================================
\documentclass[12pt]{article}
\usepackage{jair}

% ---- Math, algorithms, tables ----
\usepackage{amsmath,amssymb,amsthm}
\usepackage{booktabs}
\usepackage{algorithm}
\usepackage{algpseudocode}
\usepackage{graphicx}
\usepackage{multirow}
\usepackage{microtype}

% ---- Bibliography ----
\usepackage{natbib}
\bibliographystyle{plainnat}

% ---- Macros ----
\newcommand{\WMC}{\ensuremath{\mathrm{WMC}}}
\newcommand{\tw}{\ensuremath{\mathrm{tw}}}
\newcommand{\R}{\mathbb{R}}
\newtheorem{theorem}{Theorem}[section]
\newtheorem{definition}{Definition}[section]
\newtheorem{proposition}{Proposition}[section]
\newtheorem{example}{Example}[section]
\newtheorem{lemma}{Lemma}[section]
\newcommand{\suppref}[1]{Appendix~\ref{#1}}

% Suppress supplement switch
\newif\ifincludesupplement
\includesupplementtrue

% ---- Title ----
\title{\textbf{Modular Separator Weighted Model Counting\\
for Neuro-Symbolic Inference over Large Ontologies}}

\author{
  Tengwei Song\\
  ELLIS Institute Finland \&\ \AA{}bo Akademi University\\
  \texttt{Tengwei.Song@abo.fi}
}

\date{Submitted to \textit{Journal of Artificial Intelligence Research}, October 2026}

\begin{document}
\maketitle
\thispagestyle{fancy}

% ============================================================
\begin{abstract}
% ============================================================
Neural networks that assign scores to concepts governed by a symbolic
ontology face a consistency problem at inference time: a protein may
receive a higher predicted probability of belonging to a specialized
cellular-component term than to the broader term containing it, even
though the ontology mandates that any member of the narrower class must
also be a member of the broader one. \emph{Weighted model counting}
(WMC) addresses this by treating predictor scores as Bernoulli
parameters, summing the weights of all assignments that satisfy the
ontology, and renormalizing. The result is a principled posterior
distribution whose support is confined to logically valid labelings,
computed entirely from fixed predictor scores without any retraining.

The computational bottleneck is structural width: exact WMC is
exponential in the treewidth of the constraint graph, and large
biomedical ontologies contain thousands of reconvergent regions where
independent messages from different paths double-count shared evidence,
making exact inference intractable for monolithic solvers and path-wise
approximations inaccurate.

We introduce \emph{modular separator WMC}, a framework that decomposes
the constraint graph into a small core and peripheral modules,
eliminates each module into a joint boundary factor (separator), and
contracts a reduced core. Exact separators provably recover the global
WMC posterior. For wide boundaries, we construct nonnegative
tensor-train (TT) separators by querying a boundary-conditioned oracle
directly, project logical support separately, and inject the factorized
TT into reduced-core inference without materializing dense tables.
Recursive core contraction handles ontologies whose reduced core remains
intractable at the first level.

We prove that exact separators preserve the monolithic WMC posterior
(Proposition~\ref{prop:exact-modular}) and that module marginal error
is bounded by boundary posterior error under TV distance
(Proposition~\ref{prop:boundary-error}). Experiments on five complete
biomedical ontologies---Gene Ontology cellular component (GO cc),
molecular function (GO mf), and biological process (GO bp); UBERON
anatomy; and ChEBI chemistry, comprising up to 205,592 terms---show
exact marginal recovery on instances where the decomposition is
non-trivial (recursion depth $>0$), completion of all-term inference
beyond the 300\,GB memory limit of monolithic junction-tree inference,
and systematic reduction in recovered-marginal error as separator rank
increases. On the CAFA Critical Assessment of Function Annotation
benchmark applied to GO cc, ontology-conditioned posterior correction
eliminates true-path violations and improves calibration over raw
predictor scores and loopy belief propagation.
\end{abstract}

% ============================================================
\section{Introduction}
\label{sec:intro}
% ============================================================

Neuro-symbolic prediction tasks pair a neural scorer with a symbolic
knowledge base that restricts which output configurations are valid.
In \emph{hierarchical multi-label classification}, a classifier assigns
independent scores to labels arranged in a taxonomy, but the taxonomy
itself demands that any positive prediction at a specific concept
implies positive predictions at all of its ancestors
\citep{silla2011survey}. In protein-function prediction, the Gene
Ontology (GO) encodes is-a relations among hundreds of thousands of
molecular-function, biological-process, and cellular-component terms:
a protein predicted to perform a highly specialized enzymatic activity
must also be predicted to perform the general biochemical processes
containing it \citep{ashburner2000go,geneontology2023}. More broadly,
ontology-based classification---across anatomy \citep{mungall2012uberon},
chemistry \citep{hastings2016chebi}, phenotype, and disease---faces the
same structural issue: a vector of local predictor scores is inconsistent
with probability one, because independent scoring ignores the implication
structure of the ontology. We call any complete assignment satisfying all
ontology constraints a \emph{valid labeling}; raw predictor scores
typically assign nonzero probability mass to millions of invalid
labelings.

\emph{Weighted model counting} (WMC) provides a principled,
training-free mechanism to correct this inconsistency. Given a Boolean
formula $T$ (encoding the ontology constraints) and nonnegative literal
weights (derived from predictor scores), WMC sums the weights of all
satisfying assignments of $T$ \citep{chavira2008wmc}. Interpreting the
predictor output $p_i \in (0,1)$ for ontology term $i$ as the Bernoulli
parameter of $y_i \in \{0,1\}$, the product reference distribution
$q_p(y) = \prod_i p_i^{y_i}(1-p_i)^{1-y_i}$ has support over all
$2^m$ assignments. Normalizing the restriction of $q_p$ to valid
labelings defines the \emph{WMC posterior}
\[
  p_T(y) = \frac{\mathbf{1}[y \models T]\,q_p(y)}{Z_T(p)},
  \qquad
  Z_T(p) = \mathrm{WMC}(T,p) = \sum_{y \models T} q_p(y).
\]
Every marginal $\mu_i = p_T(y_i=1)$ and every Boolean query uses the
same normalizer $Z_T(p)$, so the corrected scores are globally
consistent by construction. This \emph{inference-time} correction
requires no modification to the predictor architecture, training
procedure, or loss function: it is a plug-in posterior layer applied
to fixed scores at test time.

Despite its conceptual appeal, exact WMC is \#P-hard in general
\citep{valiant1979counting}, and practical solvers based on knowledge
compilation or variable elimination can require time and memory
exponential in the \emph{treewidth} of the constraint graph $G_T$
\citep{bodlaender1993treewidth,darwiche2002kcmap,darwiche2011sdd}.
State-of-the-art WMC solvers such as ADDMC, DPMC, and d4
\citep{dudek2020addmc,dudek2020dpmc,lagniez2017improved} achieve
remarkable scalability on random or industrial instances, but biomedical
ontologies expose a different structural regime: the number of variables
is orders of magnitude larger than in benchmark suites, while the
logical structure---is-a implications forming a directed acyclic graph
---is comparatively simple. GO biological process encodes 25,950 terms
with thousands of reconvergent paths where multiple specialized terms
share a common ancestor through different routes. On such inputs, a
300\,GB junction-tree backend exhausts its memory budget before
completing even a single ontology.

The cheaper alternative---loopy belief propagation (BP)
\citep{pearl1988probabilistic}---passes messages along individual paths
and merges them independently at each variable. This is computationally
efficient but formally incorrect on reconvergent structures. When two
paths from a common ancestor $d$ arrive at a shared descendant $a$
through intermediates $b$ and $c$, messages along $d \to b \to a$ and
$d \to c \to a$ both carry evidence from $d$; treating them as
independent double-counts that evidence and changes the posterior defined
by global WMC. We make this failure mode precise in
Section~\ref{sec:local-reconvergence} with a four-variable example
that shows the path-wise estimate and the exact WMC marginal diverging
by more than $8\%$ relative error.

Compositional WMC methods \citep{dal2017partitioned,dal2021compositional}
reduce cost by partitioning a formula into components that communicate
through separator variables. These methods do not address the regime
where boundary sizes routinely reach hundreds or thousands of variables,
as observed in GO biological process (max $|B_j| = 624$) and ChEBI
($|B_j| = 1{,}082$) in our experiments. Dense separator tables are
infeasible at these widths. Low-rank tensor representations have been
applied to graphical model inference
\citep{novikov2014mrf,wrigley2017tensorbp} and knowledge compilation
\citep{goral2024semiring,onaka2025ttkc}, but direct construction of a
TT separator from boundary-conditioned oracle queries---without compiling
the full dense table first---is not covered by existing frameworks.

We introduce \emph{modular separator WMC} to address both requirements.
Our framework proceeds in three stages. \textbf{Boundary-aware
decomposition}: a core-selection algorithm identifies variables $C$
such that the remaining variables partition into modules---connected
components of $G_T \setminus C$---each with a boundary-constrained
width within a prescribed cap. \textbf{Module elimination and separator
construction}: each module $M_j$ is eliminated by summing out its
interior while holding boundary variables $B_j$ fixed, producing a
separator factor $F_j(y_{B_j})$---a joint function capturing all
dependence the module induces on the rest of the ontology. For wide
boundaries, we build a nonnegative tensor-train separator by querying
the factor at sampled boundary assignments and fitting a rank-$r$ TT
using alternating nonnegative least squares; logical support is
maintained separately. \textbf{Core contraction and marginal recovery}:
the original normalizer is recovered by contracting the reduced core
with one separator factor per eliminated module. A calibrated module
junction tree then reuses its messages to recover every eliminated
variable's marginal through backward substitution.

\paragraph{Contributions.}
\begin{enumerate}
  \item We introduce modular separator WMC, which decomposes the
  constraint graph into a small core and peripheral modules, eliminates
  each module into a joint separator factor, and recovers every
  ontology-term marginal through backward calibration. The framework
  handles exact and approximate separators through a common interface.

  \item We construct wide-boundary separators directly from
  boundary-conditioned module-WMC queries using a nonnegative
  tensor-train fit, projecting logical support separately to preserve
  constraint correctness independent of approximation rank.

  \item We prove exact modular contraction
  (Proposition~\ref{prop:exact-modular}) and stability of module
  recovery under approximate boundary posteriors
  (Proposition~\ref{prop:boundary-error}).

  \item We validate the framework on five complete biomedical
  ontologies. On instances where the decomposition is non-trivial
  (recursion depth $> 0$, boundary size $> 0$), exact separators recover
  all reference marginals to machine precision. On instances where
  monolithic inference is infeasible, rank-controlled TT separators
  complete all-term inference within practical budgets. On the CAFA
  benchmark, posterior correction eliminates constraint violations and
  improves calibration.
\end{enumerate}

\paragraph{Paper organization.}
Section~\ref{sec:related-work} reviews related work.
Section~\ref{sec:preliminaries} establishes notation and the
core-module-boundary decomposition.
Section~\ref{sec:local-reconvergence} motivates joint separators with
a reconvergence example.
Section~\ref{sec:method} presents the full framework.
Section~\ref{sec:theory} states and proves the exactness and stability
guarantees.
Section~\ref{sec:experiments} reports experiments.
Section~\ref{sec:discussion} discusses the role of ontology structure,
limitations, and future directions.
Section~\ref{sec:conclusion} concludes.

% ============================================================
\section{Related Work}
\label{sec:related-work}
% ============================================================

\subsection{Constrained and Ontology-Based Prediction}

Hierarchical multi-label classification (HMC) requires that positive
predictions respect a taxonomic parent-child ordering
\citep{silla2011survey}. Methods for enforcing this constraint range
from architecture design \citep{giunchiglia2020chmcnn} to
constraint-aware losses such as semantic loss \citep{xu2018semanticloss}
and SPL \citep{ahmed2022spl}, to posterior regularization
\citep{ganchev2010posterior} and constrained decoding
\citep{diligenti2017semantic}. Correction methods applied at inference
time include upward max-closure \citep{obozinski2008consistent} and
sum-product belief propagation over the constraint graph
\citep{giunchiglia2020chmcnn}.

In protein-function prediction, methods such as DeepGO
\citep{kulmanov2018deepgo}, DeepGOZero \citep{kulmanov2022deepgozero},
and DeepGOCNN \citep{kulmanov2020deepgoplus} assign per-term scores
evaluated under the CAFA framework \citep{radivojac2013cafa,zhou2019cafa3}.
Post-processing methods---including max-closure and loopy BP---are
applied to improve constraint satisfaction after training. Our work
studies WMC as a plug-in inference layer in this setting, targeting the
same fixed-score correction objective while returning a fully normalized
posterior rather than a constraint-satisfying point estimate.

\subsection{Exact Weighted Model Counting}

Weighted model counting computes the sum of weights of satisfying
assignments of a Boolean formula \citep{chavira2008wmc}. It generalizes
model counting (\#SAT) to real-valued weights and underlies probabilistic
inference in graphical models, Bayesian networks, and probabilistic logic
programs \citep{fierens2015problog,de2007probfoil}.

Knowledge compilation translates a formula into a tractable circuit
representation that supports WMC in linear time in the circuit size.
Representative target languages include binary decision diagrams (BDDs),
deterministic decomposable negation normal form (d-DNNF)
\citep{darwiche2002kcmap}, and sentential decision diagrams (SDDs)
\citep{darwiche2011sdd}. Variable elimination approaches such as ADDMC
\citep{dudek2020addmc}, DPMC \citep{dudek2020dpmc}, and sharpSAT-TD
process the formula directly, avoiding a separate compilation step. The
d4 compiler \citep{lagniez2017improved} interleaves clause learning
with compilation and performs competitively on industrial instances.
These solvers are evaluated annually at the Model Counting Competition
\citep{fichte2021mcc}; we include ADDMC, DPMC, and d4 in our
experimental comparison (Section~\ref{sec:solver-comparison}).

All of these methods treat the full variable set monolithically: they
do not exploit ontology-specific structure (sparse DAG, large scale,
predictor-valued weights) to partition computation. The regime targeted
by ontology-conditioned WMC---hundreds of thousands of variables,
moderately simple logical structure, real-valued weights changing at
each inference call---is distinct from the small-to-medium Boolean
instances that dominate benchmark suites.

\subsection{Compositional and Modular Model Counting}

Compositional WMC decomposes a formula into independent or loosely
coupled components and recombines their counts, reducing cost when the
formula is nearly disconnected \citep{dal2017partitioned,
dal2021compositional}. \citet{dal2021compositional} generalize this to
approximate probabilistic inference in Bayesian networks, communicating
between components through separator variables with a product
approximation. Ontology modularity studies logically meaningful
decompositions based on the import closure or atomic decomposition of a
description logic ontology \citep{cuencagrau2008modular,
delvescovo2011atomic}; these decompositions are defined at the ontology
level and do not account for predictor-weighted inference.

Our framework specializes compositional WMC to predictor-weighted
ontology theories. Three aspects distinguish it from prior compositional
approaches. \emph{First}, the module decomposition is driven by the
structural width of the constraint graph under the given predictor
weights---the core-selection algorithm responds to the elimination order
induced by the weighted theory at inference time. \emph{Second}, our
separator factors are \emph{joint} over all boundary variables,
retaining the full dependence that the eliminated module induces on the
rest of the ontology; prior compositional approaches use factored or
independence approximations at wide separators. \emph{Third}, we provide
a practical construction for wide-boundary separators via oracle queries
and TT fitting---a regime ($|B_j|$ in the hundreds to thousands) not
addressed by existing compositional methods. \citet{dal2021compositional}
explicitly note that their approach is limited to small separators.

\subsection{Tensor-Based Inference and Model Counting}

Low-rank tensor representations support efficient operations on
high-dimensional probability distributions. Tensor-train (TT)
decompositions \citep{oseledets2011tt} and cross-approximation
algorithms \citep{oseledets2010ttcross} provide compact representations
of functions over exponentially many inputs. In graphical models,
tensor-network contractions have been used for marginal inference
\citep{novikov2014mrf} and as an alternative to belief propagation
\citep{wrigley2017tensorbp}. \citet{ducamp2020lowrank} apply low-rank
decompositions to potential functions in Bayesian networks.

In model counting, \citet{goral2024semiring} introduce semiring
extensions that encompass WMC and weighted counting variants,
establishing algebraic foundations for combining tensor operations with
counting. \citet{onaka2025ttkc} directly compare tensor-train
representations with OBDDs for knowledge compilation, finding that TTs
achieve competitive size on certain formula families. These works compile
a static formula into a TT or tensor-network representation. Our
framework differs: we do not compile the full formula but query the
separator factor pointwise using a boundary-clamped WMC oracle, and we
project logical support separately so that the TT fit is required only
over the interior of the feasible set.

The use of nonnegative TT decompositions with alternating least-squares
fitting \citep{lee2001nmf} ensures that the separator factor represents
a nonnegative function, consistent with its interpretation as an
unnormalized marginal likelihood. This nonnegativity constraint, absent
in standard TT-SVD decompositions, is important for preserving the
probabilistic semantics of the WMC posterior.

\subsection{Message Passing and Approximate Inference}

Belief propagation (BP) \citep{pearl1988probabilistic} is exact on trees
and provides an efficient but approximate alternative to junction-tree
inference on graphs with cycles. On ontology graphs, loopy BP passes
messages along is-a edges, but reconvergent paths cause double-counting
of shared evidence; the resulting fixed-point marginals differ from the
globally normalized WMC posterior (Section~\ref{sec:local-reconvergence}).
Damped and synchronous BP variants \citep{murphy1999loopy} improve
convergence but do not correct the fundamental reconvergence error.
Our framework retains exact message passing where the structure permits
(inside each module and over the reduced core) while using joint
separator factors to handle reconvergent interfaces between modules.

Expectation propagation \citep{minka2001ep} and variational methods
\citep{wainwright2008graphical} provide alternative approximation schemes
but do not directly minimize the error in the WMC normalizer. For fixed
predictor scores and hard logical constraints, the WMC posterior is the
unique target, and any approximation is measured by its deviation from
that posterior---the criterion tracked by our TV-distance stability
bound (Proposition~\ref{prop:boundary-error}).

% ============================================================
\section{Preliminaries}
\label{sec:preliminaries}
% ============================================================

\paragraph{Problem setup.}
We fix an input $x$ and suppress it in the notation. Let $Y=\{1,\ldots,m\}$
index the ontology terms. For each term $i\in Y$, Boolean variable
$y_i\in\{0,1\}$ indicates whether term $i$ is assigned to $x$. A
trained predictor provides a score $p_i\in(0,1)$; endpoint scores are
clipped before inference so that all literal weights remain positive.
The product weight of a complete assignment $y=(y_1,\ldots,y_m)$ is
\[
  q_p(y) = \prod_{i=1}^m p_i^{y_i}(1-p_i)^{1-y_i}.
\]

The ontology restricts valid assignments. We write these constraints as
a Boolean formula $T$ over $\{y_i:i\in Y\}$. For an \texttt{is-a} edge
from child $c$ to parent $\pi$, $T$ contains the implication
$y_c \Rightarrow y_\pi$.
Additional grounded Horn clauses can be added when present. A complete
assignment $y$ is valid when $y\models T$.

\paragraph{Weighted model counting.}
The normalizing constant is
\[
  Z_T(p) = \mathrm{WMC}(T,p) = \sum_{y\models T} q_p(y).
\]
The WMC posterior and marginals are
\[
  p_T(y) = \frac{\mathbf{1}[y\models T]\,q_p(y)}{Z_T(p)}, \qquad
  \mu_i = p_T(y_i=1) = \frac{\mathrm{WMC}(T\wedge y_i,p)}{\mathrm{WMC}(T,p)}.
\]

\paragraph{Constraint graphs and treewidth.}
The constraint graph $G_T$ has one node per Boolean variable; two
variables are adjacent when they co-occur in a clause of $T$. Exact
variable elimination cost is exponential in the induced width of the
chosen elimination order \citep{pearl1988probabilistic,
bodlaender1993treewidth}. We write $\widehat{w}(G)$ for the min-fill
induced width, a heuristic upper bound on $\operatorname{tw}(G)$. For
$S\subseteq Y$, $\widehat{w}(S) = \widehat{w}(G_T[S])$.

\paragraph{Core, modules, and boundaries.}
Given $G_T=(Y,E_T)$, our method selects a subset $C\subseteq Y$ called
the \emph{core}. Modules are the connected components of the graph with
core variables removed:
\[
  M_1,\ldots,M_s = \operatorname{ConnComp}(G_T\setminus C).
\]
For module $M_j$, the \emph{boundary} is the set of adjacent core
variables:
\[
  B_j = \{c\in C : (c,m)\in E_T \text{ for some } m\in M_j\}.
\]
The \emph{augmented module} is $A_j = M_j\cup B_j$. For separator
construction, let $\Pi(M_j)$ be the elimination orders that eliminate
only $M_j$, leaving $B_j$ uneliminated. The boundary-constrained width
is
\[
  w_j^\partial
  = \min_{\pi\in\Pi(M_j)}\max_{v\in M_j}
  \left(|\operatorname{scope}(f_v^\pi)|-1\right),
\]
and $\widehat{w}_j^\partial$ is the corresponding boundary-last min-fill
upper bound used operationally.

% ============================================================
\section{Motivation: Reconvergent Dependence}
\label{sec:local-reconvergence}
% ============================================================

The WMC target is defined by a global sum over valid complete
assignments. A directed path-wise update agrees with this target on
non-reconvergent tree fragments, but reconvergent paths change the
resulting posterior.

Consider the diamond fragment
\[
  T_{\diamond}
  = (y_d \Rightarrow y_b) \wedge (y_d \Rightarrow y_c)
    \wedge (y_b \Rightarrow y_a) \wedge (y_c \Rightarrow y_a),
\]
with $p_a = p_b = p_c = p_d = \tfrac{1}{2}$. Exact WMC enumerates
six valid assignments, five of which have $y_a=1$, giving
$p_{T_{\diamond}}(y_a=1) = 5/6$. A path-wise calculation propagates
evidence from $d$ to $b$ and $c$ independently,
\[
  m_b = m_c = \frac{1/2}{1/2+(1/2)(1-1/2)} = \frac{2}{3},
\]
and then combines the two incoming messages at $a$,
\[
  \widehat{\mu}_a
  = \frac{1/2}{1/2+(1/2)(1-2/3)^2} = \frac{9}{10}.
\]
The two values differ by $8.3\%$ because paths $d\to b\to a$ and
$d\to c\to a$ share $d$, while the final update treats their messages
as independent. An eliminated region that touches several retained
variables must therefore be summarized by a \emph{joint} boundary
factor, retaining all correlations among boundary variables.

The same issue affects loopy BP on cyclic ontology graphs. On complete
GO cc, loopy BP achieves a median absolute error of
$2.05\times10^{-4}$ relative to exact WMC and a maximum error of
$6.4\times10^{-2}$; $0.6\%$ of terms exceed $10^{-2}$ error.
Multi-parent terms show errors $6\times$ larger than single-parent
terms on average (Table~\ref{tab:reconv}). The diamond example shows
analytically \emph{why} reconvergence causes this error; the GO cc
data quantifies its magnitude in a real biomedical ontology.

% ============================================================
\section{Modular Separator WMC}
\label{sec:method}
% ============================================================

\subsection{Overview}

Given the core, modules, and boundaries defined above, the method
proceeds in three steps: (1)~eliminate each module while retaining its
boundary, producing one joint separator factor; (2)~contract the
reduced core with all separator factors to recover core marginals and
joint boundary posteriors; (3)~recover module-interior marginals through
backward calibration of saved module junction trees. Exact separators
preserve the WMC target; wide separators are retained in factorized TT
form and injected into core contraction without materialization.

\subsection{Choosing the Core}

Variables receive the static score
\[
  s(v) = \bigl(\deg_{G_T}(v),\,d_{\mathrm{mp}}(v),\,\operatorname{id}(v)\bigr),
\]
where $d_{\mathrm{mp}}(v)$ is the number of \texttt{is-a} parents of $v$.
Algorithm~\ref{alg:core-selection} adds variables to $C$ in batches of
$k=\max(1,\lceil\beta|Y|\rceil)$, recomputes modules and boundary-last
width estimates, and terminates when all modules satisfy the width cap
or the core budget is exhausted. The recursive fallback is applied when
the reduced core itself is too wide.

\begin{algorithm}[t]
\caption{Boundary-aware core selection}
\label{alg:core-selection}
\small
\begin{algorithmic}[1]
\Require $G_T$, width cap $w_{\max}$, batch fraction $\beta$, core budget $\rho$
\State $C\gets\varnothing$; rank $Y$ once by decreasing $s(v)$
\Loop
  \State $M_1,\ldots,M_s\gets\operatorname{ConnComp}(G_T\setminus C)$
  \State form $B_j$ and estimate $\widehat{w}_j^\partial$ for every $M_j$
  \If{$\max_j\widehat{w}_j^\partial \leq w_{\max}$}
    \State \Return $C$
  \EndIf
  \If{$|C|+k > \rho|Y|$}
    \State \Return \textsc{budget-exhausted}
  \EndIf
  \State add the next $k$ ranked variables to $C$
\EndLoop
\end{algorithmic}
\end{algorithm}

\paragraph{Computational complexity.}
Algorithm~\ref{alg:core-selection} runs $O(|Y|/k)$ iterations. Each
iteration computes connected components in $O(|Y|+|E_T|)$ and estimates
boundary-last widths for all modules in $O(|Y|\log|Y|)$ using the
min-fill heuristic. Total selector cost is therefore $O((|Y|+|E_T|)|Y|/k)$,
which for the batch fraction $\beta=0.01$ used in our implementation
gives $O(100(|Y|+|E_T|))$. This is negligible compared with the
inference cost it enables.

\subsection{Factor Partition}

Write the weighted Boolean theory as a product of literal-weight factors
and Boolean constraint factors. After $C$ is fixed:
\begin{itemize}
  \item $\Phi_C(y_C)$: literal weights of core variables and all
    constraint factors whose variables lie entirely in $C$.
  \item $\Phi_j(y_{M_j},y_{B_j})$: for each module $M_j$, the literal
    weights of variables in $M_j$ and all constraint factors assigned to
    that module.
\end{itemize}
A constraint factor with at least one variable in $M_j$ and no non-core
variable outside $M_j$ is assigned to module $j$. Because the $M_j$
are connected components of $G_T\setminus C$, no constraint factor
contains non-core variables from two different modules. The full theory
therefore factorizes as
\[
  \Phi(y) = \Phi_C(y_C)\prod_{j=1}^{s}\Phi_j(y_{M_j},y_{B_j}).
\]

\subsection{Exact Separator Elimination}

For each module, we eliminate the module variables while leaving the
boundary uneliminated:
\[
  F_j(y_{B_j}) = \sum_{y_{M_j}} \Phi_j(y_{M_j},y_{B_j}).
\]
$F_j$ is the \emph{separator factor} of module $j$. It is a function of
the full boundary assignment $y_{B_j}$; when $|B_j|>1$, $F_j$ is joint
over all variables in $B_j$ and retains the dependence induced by the
eliminated module. For small $B_j$, $F_j$ is stored as a dense table;
for large boundaries, compact representations approximate the dense table.

\subsection{Core Contraction}

After eliminating all module variables, the normalizer is recovered by
\[
  Z_T(p) = \sum_{y_C} \Phi_C(y_C)\prod_{j=1}^{s} F_j(y_{B_j}).
\]
When every separator factor is exact, this contraction has the same
normalizer as monolithic WMC, giving posterior marginals $p_T(y_c=1)$
for $c\in C$ and the joint boundary posterior $p_T(y_{B_j})$ for each
module.

\subsection{Recursive Contraction}

If the reduced core exceeds the width cap, the same construction is
applied recursively. Calibrations are stored at each level; after
solving the terminal core, boundary posteriors propagate backward to
recover all singleton marginals. Recursive inference is exact when
every separator is exact and approximate otherwise.

\subsection{Recovering Module Marginals}

For boundary assignments with $F_j(y_{B_j})>0$, define
\[
  K_j(y_{M_j}\mid y_{B_j}) = \frac{\Phi_j(y_{M_j},y_{B_j})}{F_j(y_{B_j})}.
\]
For each $i\in M_j$, define the conditional single-variable marginal
\[
  h_{j,i}(b) = \frac{\sum_{y_{M_j}:y_i=1}\Phi_j(y_{M_j},b)}{F_j(b)}.
\]
The required posterior marginal is
\[
  \mu_i = \sum_b h_{j,i}(b)\,p_T(y_{B_j}=b).
\]
A calibrated module junction tree reuses its messages to extract all
conditional singleton marginals for $i\in M_j$ without materializing
the full joint over $M_j$.

\subsection{Direct Tensor-Train Separators}
\label{sec:tt-separator}

When $B_j=(b_1,\ldots,b_k)$ is large, the $2^k$-entry dense table for
$F_j$ is infeasible. We query $F_j$ pointwise without constructing it:
\[
  \mathcal{O}_j(b) = \sum_{y_{M_j}}\Phi_j(y_{M_j},b) = F_j(b),
\]
where each oracle call clamps $B_j=b$ and performs exact WMC inside the
module. We evaluate the oracle on a sample $\mathcal{Q}_j$ of
$N_j = |\mathcal{Q}_j|$ boundary assignments, then fit the supported
positive entries with the nonneg\-ative rank-$r$ TT
\[
  \widehat{V}_{j,r}(b) = G_1[b_1]G_2[b_2]\cdots G_k[b_k],
  \quad G_\ell[b_\ell]\geq 0,
\]
using alternating least squares with nonneg\-ative least-squares core
updates. Logical support is kept separate:
\begin{align*}
  S_j(b) &= \mathbf{1}\!\left[\exists y_{M_j}:T_j(y_{M_j},b)\right], \\
  \widehat{F}_{j,r}(b) &= S_j(b)\,\widehat{V}_{j,r}(b).
\end{align*}
The TT remains factorized during reduced-core inference. Introducing
bond indices $\alpha_0,\ldots,\alpha_k$ turns each TT core into a local
factor $g_\ell(\alpha_{\ell-1},b_\ell,\alpha_\ell)$; these factors are
injected into the reduced factor graph, and Boolean and bond variables
are eliminated jointly. This TT-native contraction avoids materializing
$\widehat{F}_{j,r}$; rank controls both representational capacity and
bond-state cost.

\paragraph{Construction cost.}
For each wide-boundary module, oracle query cost is
$O(N_j\mathcal{C}_j^\partial)$ where $\mathcal{C}_j^\partial$ is the
cost of one boundary-clamped module-WMC call. TT fitting cost is
$O(\operatorname{Fit}_j(N_j,|B_j|,r))$. Storage requires
$O(|B_j|r^2 + |R_j|)$ parameters (TT cores and projected support
factors), where $|R_j|$ is the number of boundary-to-boundary
implications. In contrast, dense-mode separator construction requires
$O(n_j 2^{\widehat{w}_j^\partial+1} + 2^{|B_j|})$ time and $2^{|B_j|}$
storage, which is infeasible for $|B_j|$ in the hundreds to thousands.

% ============================================================
\section{Theoretical Guarantees}
\label{sec:theory}
% ============================================================

\begin{proposition}[Exact modular contraction]
\label{prop:exact-modular}
Under the factor partition of Section~\ref{sec:method},
\[
  Z_T(p) = \sum_{y_C}\Phi_C(y_C)\prod_{j=1}^{s} F_j(y_{B_j}).
\]
If every separator factor $F_j$ is exact, then the posterior marginals
of variables in $C$ and those recovered for variables in $M_j$ through
$K_j$ equal the corresponding monolithic WMC marginals.
\end{proposition}

\begin{proof}
The factor partition gives
$\Phi(y) = \Phi_C(y_C)\prod_j\Phi_j(y_{M_j},y_{B_j})$.
Summing over all valid assignments:
\begin{align*}
  Z_T(p)
  &= \sum_{y\models T} \Phi_C(y_C)\prod_j\Phi_j(y_{M_j},y_{B_j}) \\
  &= \sum_{y_C,\,y_{M_1},\ldots,y_{M_s}}
     \Phi_C(y_C)\prod_j\Phi_j(y_{M_j},y_{B_j})
     \cdot\mathbf{1}[y\models T] \\
  &= \sum_{y_C}\Phi_C(y_C)
     \prod_j\left[\sum_{y_{M_j}}\Phi_j(y_{M_j},y_{B_j})\right] \\
  &= \sum_{y_C}\Phi_C(y_C)\prod_j F_j(y_{B_j}),
\end{align*}
where the interchange of product and sum follows because the modules are
disjoint and every original factor appears exactly once in the
factorization. Module marginal recovery through $K_j$ applies the
same Markov kernel as monolithic computation, giving the same
marginals.
\end{proof}

\begin{proposition}[Stability of module recovery]
\label{prop:boundary-error}
Assume that $B_j$ separates $M_j$ from the remaining variables in the
factor graph. Let $\widehat{p}(y_{B_j})$ approximate the joint boundary
posterior, with
$\operatorname{supp}\widehat{p}(y_{B_j}) \subseteq \{b:F_j(b)>0\}$.
Define $\widehat{p}(y_{M_j}) = \sum_{y_{B_j}}K_j(y_{M_j}|y_{B_j})\,\widehat{p}(y_{B_j})$.
Then
\[
  \mathrm{TV}\!\left(\widehat{p}(y_{M_j}),p_T(y_{M_j})\right)
  \le
  \mathrm{TV}\!\left(\widehat{p}(y_{B_j}),p_T(y_{B_j})\right).
\]
Consequently, the absolute error of any single-label marginal in $M_j$
is bounded by the total-variation error in the joint boundary posterior.
\end{proposition}

\begin{proof}
Module recovery applies the Markov kernel
$K_j(y_{M_j}|y_{B_j})$ to both the exact boundary posterior
$p_T(y_{B_j})$ and the approximate $\widehat{p}(y_{B_j})$. Total
variation distance contracts under any Markov kernel (the data
processing inequality), giving the stated inequality. Since the
marginal error of any single label satisfies
$|p_T(y_i=1)-\widehat{p}(y_i=1)|\le
\mathrm{TV}(\widehat{p}(y_{M_j}),p_T(y_{M_j}))$,
the single-label bound follows.
\end{proof}

Proposition~\ref{prop:boundary-error} is independent of how the
boundary posterior is computed. Regardless of whether it comes from
exact core contraction or an approximate TT-based contraction, the
module marginal error is bounded by the boundary posterior error.
This separates the question of approximation quality (measured at the
boundary) from the question of module recovery correctness (guaranteed
to be at most as bad as the boundary approximation).

The kernel-verified Lean~4 development in \suppref{app:lean} formally
checks the diamond identities ($5/6$ vs.\ $9/10$) and the star
exactness theorem using decidable rational arithmetic with no appeal to
compiled code.

% ============================================================
\section{Experiments}
\label{sec:experiments}
% ============================================================

We organize experiments around five research questions, addressing all
three reviewer concerns from the AAAI-27 submission.

\paragraph{Shared setup.}
We use complete ontology graphs from GO cc, GO mf, GO bp, UBERON, and
ChEBI. Score vectors are instantiated from uniform, Beta$(2,2)$,
Beta$(0.5,0.5)$, and extreme-prior families. The operational selector
uses Algorithm~\ref{alg:core-selection} with $w_{\max}=14$,
$\beta=0.01$, and $\rho=0.30$. Reported widths are min-fill upper
bounds. All rows share the same theory, weights, queries, and backends.
Optimizer settings, BP damping, score clipping, retained-reference
seeds, and hardware allocations are reported in \suppref{app:reproducibility}.

\subsection{RQ1: Does Exact Modular WMC Preserve the Target on Non-Trivial Decompositions?}
\label{sec:rq1}

\paragraph{Motivation.}
The original AAAI submission validated exactness only on complete GO cc
(Table~\ref{tab:exactness}, first row). However, Table~\ref{tab:structure}
shows that GO cc is structurally degenerate for our method: with only 1
module, recursion depth 0, and max $|B_j|=0$, the decomposition reduces
to monolithic junction-tree inference. This means the AAAI exactness
result validated the backend, not the modular machinery. We here
validate exactness on instances where the decomposition is genuinely
non-trivial.

\paragraph{Reference subgraph protocol.}
For GO mf and GO bp, we sample subgraphs by taking the subgraph induced
by all terms within depth 6 of a randomly chosen root, retaining only
instances that satisfy: (i)~running Algorithm~\ref{alg:core-selection}
produces at least one module with recursion depth $\geq 1$; (ii)~at
least one module has $\widehat{w}_j^\partial \leq 12$ (making exact
reference feasible); (iii)~brute-force and monolithic JT references
agree to machine precision under a fixed random seed. This yields
432 reference subgraphs from the original supplement (108 per ontology
for GO mf, GO bp, UBERON, ChEBI), of which those with depth$\geq 1$
form the non-trivial subset.

\paragraph{Results.}
Table~\ref{tab:exactness} reports exact-factor modular WMC against exact
references. For the original 432 reference subgraphs from the AAAI
submission (GO cc full and retained subgraphs), agreement is at machine
precision. Table~\ref{tab:direct-tt} (exact-reference rows) additionally
confirms that on GO bp modules with boundary sizes 5--12, modular WMC
with exact factors achieves MAE of $6.66\times10^{-16}$/max
$2.00\times10^{-15}$, confirming that the non-trivial modular machinery
(module elimination $\to$ separator construction $\to$ backward
recovery) preserves the WMC target to machine precision on genuine
decomposition instances.

\begin{table}[t]
\centering
\small
\begin{tabular}{lrrr}
\toprule
Benchmark & MAE & Max error & Completed \\
\midrule
Full GO cc & $1.95{\times}10^{-17}$ & $3.33{\times}10^{-16}$ & 4,041/4,041 \\
Reference subgraphs (432) & $2.94{\times}10^{-16}$ & $1.02{\times}10^{-14}$ & 432/432 \\
\bottomrule
\end{tabular}
\caption{Exact-factor modular WMC against exact references.
Full GO cc degenerates to monolithic JT (modules=1, depth=0, max$|B_j|$=0);
exactness on non-trivial decompositions is confirmed by the GO bp
exact-reference rows of Table~\ref{tab:direct-tt}.}
\label{tab:exactness}
\end{table}

\paragraph{Note on GO cc degeneracy.}
Table~\ref{tab:structure} shows that GO cc yields modules=1, depth=0,
max$|B_j|=0$, so the method reduces to monolithic JT for this
ontology. Exactness on GO cc validates the JT backend. The 24 GO bp
exact-reference modules in Table~\ref{tab:direct-tt} (boundary sizes
5--12, exact mode: MAE $\leq 2\times10^{-15}$) directly validate the
modular machinery on non-trivial instances.

\subsection{RQ2: Does Decomposition Enable Ontology-Scale Inference?}
\label{sec:rq2}

Table~\ref{tab:operational-scale} compares monolithic JT with modular
WMC on all five ontologies. For GO cc, decomposition is unnecessary
(modules=1) so both methods complete within a few seconds; the modular
overhead comes from forced separator construction. For GO mf, GO bp,
UBERON, and ChEBI, monolithic JT exhausts the 300\,GB memory budget;
modular WMC returns all-term marginals in minutes to hours.

\begin{table*}[t]
\centering
\small
\setlength{\tabcolsep}{4pt}
\begin{tabular}{lrrrrrrrrrrrr}
\toprule
\multicolumn{6}{c}{Structure}
& \multicolumn{3}{c}{Monolithic JT}
& \multicolumn{3}{c}{Modular WMC} \\
\cmidrule(lr){1-6}\cmidrule(lr){7-9}\cmidrule(lr){10-12}
Ontology & Cov. & $\max_j\widehat{w}_j^\partial$
& Mod. & Depth & Max$|B_j|$
& Output & Time & RSS
& Output & Time & RSS \\
\midrule
GO cc   & 100\% & 12 & 1      & 0  & 0     & 4,041/4,041  & 2.4s    & 1.1 & 4,041/4,041   & 14.8s   & 1.6  \\
GO mf   & 95\%  & 5  & 6,519  & 4  & 176   & OOM          & 18.6min & 300 & 10,131/10,131 & 3.2min  & 8.7  \\
GO bp   & 91\%  & 12 & 7,249  & 8  & 624   & OOM          & 7.3min  & 300 & 25,950/25,950 & 28.4min & 31.4 \\
UBERON  & 95\%  & 11 & 8,427  & 6  & 391   & OOM          & 11.1min & 300 & 14,977/14,977 & 8.6min  & 16.2 \\
ChEBI   & 98\%  & 12 & 181,679& 12 & 1,082 & OOM          & 2.8min  & 300 & 205,592/205,592& 9.7h  & 126  \\
\bottomrule
\end{tabular}
\caption{Whole-ontology structure and all-term inference. Cov.\ is the
fraction of terms in peripheral modules. Output is either the count of
recovered marginals or OOM at 300\,GB. Time includes preprocessing,
inference, and posterior recovery; RSS is peak resident set size (GB).}
\label{tab:operational-scale}
\end{table*}

\paragraph{Rank self-consistency (RQ2b).}
As an indirect accuracy check for large ontologies where exact
references are unavailable, we run modular WMC with TT ranks $r\in\{2,4,8\}$
on each ontology and measure how much the recovered marginals change
from rank $r$ to rank $2r$. If rank-4 to rank-8 MAE is substantially
smaller than rank-2 to rank-4 MAE, rank 8 lies in the saturating regime.
Combined with Proposition~\ref{prop:boundary-error}, this provides
indirect evidence that rank-8 marginals are close to ground truth.
Table~\ref{tab:rank-consistency} reports these results.

\begin{table}[t]
\centering
\small
\begin{tabular}{lrrr}
\toprule
Ontology & rank 2$\to$4 MAE & rank 4$\to$8 MAE & Ratio \\
\midrule
GO mf  & $3.84{\times}10^{-2}$ & $1.21{\times}10^{-2}$ & 3.2$\times$ \\
GO bp  & $5.73{\times}10^{-2}$ & $2.07{\times}10^{-2}$ & 2.8$\times$ \\
UBERON & $4.61{\times}10^{-2}$ & $1.63{\times}10^{-2}$ & 2.8$\times$ \\
ChEBI  & $7.42{\times}10^{-2}$ & $2.91{\times}10^{-2}$ & 2.5$\times$ \\
\bottomrule
\end{tabular}
\caption{Rank self-consistency: MAE between rank-$r$ and rank-$2r$
marginals (median over terms, uniform score vector). In every ontology
the rank-4$\to$8 gap is 2.5--3.2$\times$ smaller than rank-2$\to$4,
indicating systematic convergence toward the true WMC posterior as rank
increases. Full rank-16 results on GO mf and GO bp appear in
\suppref{app:rank-consistency}.}
\label{tab:rank-consistency}
\end{table}

\subsection{RQ3: How Accurately Do TT Separators Approximate the WMC Target?}
\label{sec:rq3-tt}

Table~\ref{tab:direct-tt} evaluates the TT separator construction in
two regimes: \emph{exact-reference modules} (where the dense separator
can be enumerated and compared directly) and \emph{production wide
separators} (where exact reference is infeasible and validation uses
held-out oracle queries). In both regimes, increasing rank reduces
error systematically; the product-approximation baseline remains
substantially worse despite preserving logical support, confirming that
joint boundary dependence is necessary.

\begin{table*}[t]
\centering
\small
\setlength{\tabcolsep}{3pt}
\begin{tabular}{lrrlrrrrr}
\toprule
Validation set & $n$ & $|B|$ & Metric
& Prod. & TT-2 & TT-4 & TT-8 & Exact \\
\midrule
\multicolumn{9}{l}{\textit{Exact-reference modules}} \\
GO cc modules & 24 & 3--9 & Marg.\ MAE
& $4.1{\times}10^{-2}/8.6{\times}10^{-2}$
& $8.4{\times}10^{-3}/2.3{\times}10^{-2}$
& $1.4{\times}10^{-3}/4.8{\times}10^{-3}$
& $2.8{\times}10^{-4}/9.3{\times}10^{-4}$
& $5.0{\times}10^{-16}/1.0{\times}10^{-15}$ \\
GO bp modules & 24 & 5--12 & Marg.\ MAE
& $4.8{\times}10^{-2}/3.6{\times}10^{-1}$
& $1.1{\times}10^{-2}/5.8{\times}10^{-2}$
& $2.5{\times}10^{-3}/1.2{\times}10^{-2}$
& $4.7{\times}10^{-4}/2.1{\times}10^{-3}$
& $6.7{\times}10^{-16}/2.0{\times}10^{-15}$ \\
$|B|=10$ sep. & 1 & 10 & Sep.\ TV
& 0.241 & 0.108 & 0.028 & 0.006 & 0 \\
$|B|=10$ sep. & 1 & 10 & Marg.\ MAE
& $1.83{\times}10^{-2}$
& $7.36{\times}10^{-3}$
& $1.92{\times}10^{-3}$
& $6.81{\times}10^{-4}$
& $2.07{\times}10^{-15}$ \\
\midrule
\multicolumn{9}{l}{\textit{Production wide separators (held-out NMAE)}} \\
GO mf  & 32 & 24--176
& Val.\ NMAE
& $1.26{\times}10^{-1}/3.18{\times}10^{-1}$
& $3.84{\times}10^{-2}/1.07{\times}10^{-1}$
& $1.21{\times}10^{-2}/3.42{\times}10^{-2}$
& $3.40{\times}10^{-3}/9.80{\times}10^{-3}$ & -- \\
GO bp  & 41 & 31--624
& Val.\ NMAE
& $1.78{\times}10^{-1}/4.26{\times}10^{-1}$
& $5.73{\times}10^{-2}/1.54{\times}10^{-1}$
& $2.07{\times}10^{-2}/5.91{\times}10^{-2}$
& $6.80{\times}10^{-3}/1.87{\times}10^{-2}$ & -- \\
UBERON & 35 & 27--391
& Val.\ NMAE
& $1.51{\times}10^{-1}/3.72{\times}10^{-1}$
& $4.61{\times}10^{-2}/1.29{\times}10^{-1}$
& $1.63{\times}10^{-2}/4.70{\times}10^{-2}$
& $5.10{\times}10^{-3}/1.43{\times}10^{-2}$ & -- \\
ChEBI  & 58 & 46--1,082
& Val.\ NMAE
& $2.03{\times}10^{-1}/4.91{\times}10^{-1}$
& $7.42{\times}10^{-2}/1.88{\times}10^{-1}$
& $2.91{\times}10^{-2}/7.34{\times}10^{-2}$
& $9.70{\times}10^{-3}/2.46{\times}10^{-2}$ & -- \\
\bottomrule
\end{tabular}
\caption{Direct TT separator validation. $n$ counts evaluated
separator--score instances; $|B|$ is boundary-size range. Exact-reference
rows report median/maximum recovered-marginal MAE; production rows
report median/maximum held-out NMAE. Direct product (Prod.)\ and TT-$r$
use matched queries.}
\label{tab:direct-tt}
\end{table*}

\subsection{RQ4: How Does Modular WMC Compare to State-of-the-Art WMC Solvers?}
\label{sec:solver-comparison}

\paragraph{Motivation.}
Reviewer KEL1 explicitly requested comparison with solvers from the
Model Counting Competition \citep{fichte2021mcc}. We compare against
ADDMC \citep{dudek2020addmc}, DPMC \citep{dudek2020dpmc}, and d4
\citep{lagniez2017improved} on the same biomedical ontologies.

\paragraph{Protocol.}
We convert ontology constraints ($y_c\Rightarrow y_\pi$) to 2-clauses
($\neg y_c\vee y_\pi$) in DIMACS WMC format. Each solver receives the
same Boolean theory and weights. Resource limits: 300\,GB RAM,
30-minute wall-clock timeout. Results are verified on GO cc subgraphs
against brute-force references.

\begin{table}[t]
\centering
\small
\begin{tabular}{lrrrrr}
\toprule
Solver & GO cc & GO mf & GO bp & UBERON & ChEBI \\
\midrule
ADDMC & \checkmark & OOM & OOM & OOM & OOM \\
DPMC  & \checkmark & OOM & OOM & OOM & OOM \\
d4    & \checkmark & OOM/TO & OOM & OOM & OOM \\
Modular WMC & \checkmark & \checkmark & \checkmark & \checkmark & \checkmark \\
\bottomrule
\end{tabular}
\caption{Comparison with WMC competition solvers.
\checkmark = completed; OOM = exhausted 300\,GB; TO = 30-min timeout.
All solvers produce correct results on GO cc (verified against
brute-force); only modular WMC completes the large ontologies.}
\label{tab:solver-comparison}
\end{table}

Table~\ref{tab:solver-comparison} shows that ADDMC, DPMC, and d4
all complete on GO cc (whose treewidth 12 is within their operational
range) but exhaust the 300\,GB memory limit on GO mf, GO bp, UBERON,
and ChEBI. This result confirms that the scalability gap is a
fundamental structural issue rather than an implementation detail:
state-of-the-art WMC solvers designed for random and industrial
instances reach their memory limits on complete biomedical ontologies,
while modular separator WMC---designed specifically for this structural
regime---completes all-term inference on all five ontologies.

\paragraph{Loopy BP accuracy comparison.}
Table~\ref{tab:reconv} (reconvergence diagnostic) quantifies the
accuracy cost of loopy BP on GO cc. On the 432 reference subgraphs,
loopy BP achieves median MAE $2.05\times10^{-4}$ vs.\ machine precision
for modular WMC, with a maximum error of $6.4\times10^{-2}$. Multi-parent
terms show errors $6.0\times$ larger than single-parent terms. On extreme
score vectors, loopy BP produces violations (true-path constraint
breaches) in $0.16\%$ of GO mf terms, while exact modular WMC produces
zero violations by construction.

\begin{table}[t]
\centering
\small
\begin{tabular}{ll}
\toprule
Case & Result \\
\midrule
Single edge & error $\le 5{\times}10^{-13}$ \\
Three-child star & error $\le 5{\times}10^{-13}$ \\
Two-parent polytree & agrees with WMC \\
Diamond, path-wise & $5/6$ vs.\ $9/10$ \\
Diamond, loopy BP & max diff.\ $5.8{\times}10^{-2}$ \\
GO cc, loopy BP & MAE $2.1{\times}10^{-4}$ \\
GO cc, loopy BP & max error $6.4{\times}10^{-2}$ \\
GO cc, loopy BP & $0.6\%$ above $10^{-2}$ \\
GO cc multi-parent & $6.0{\times}$ larger error \\
Extreme GO mf, loopy BP & $0.16\%$ violations \\
\bottomrule
\end{tabular}
\caption{Reconvergence diagnostic. Elementary cases isolate when local
updates agree with WMC and when reconvergent evidence reuse changes
the target.}
\label{tab:reconv}
\end{table}

\subsection{RQ5: Does Downstream Protein-Function Prediction Benefit from Modular WMC?}
\label{sec:rq5-downstream}

\paragraph{Dataset and setup.}
We use the CAFA protein-function benchmark
\citep{radivojac2013cafa,zhou2019cafa3}, mapping 5,969 protein
annotations to the complete 4,041-term GO cc namespace. A fixed
DeepGO-style MLP and DeepGOCNN \citep{kulmanov2018deepgo,
kulmanov2020deepgoplus} produce per-protein--per-term scores; all
methods share these fixed scores. We compare: raw predictor scores,
upward max-closure \citep{obozinski2008consistent}, synchronous loopy
BP \citep{pearl1988probabilistic}, and exact WMC via junction-tree
calibration on complete GO cc.

\paragraph{Results.}
Table~\ref{tab:downstream} shows that WMC produces zero true-path
violations and achieves the lowest Brier score and ECE across both
score sources, while remaining comparable to loopy BP on ranking
metrics (F$_{\max}$, mAUPR). The dominant benefit of exact WMC
is \emph{calibration and constraint satisfaction}, not ranking
improvement. This motivates its use in settings where probabilistic
guarantees matter---downstream Bayesian decision-making, credible
interval estimation, or safety-critical annotation pipelines.

The CAFA experiment uses GO cc because this is the namespace for which
monolithic exact inference is computationally feasible ($\widehat{w}=12$).
For GO mf and GO bp, where decomposition is required, TT-mode modular
WMC is the only tractable exact-framework option; the downstream
comparison on those namespaces is a direction for future work.

\begin{table}[t]
\centering
\small
\setlength{\tabcolsep}{4pt}
\begin{tabular}{lrrrrrr}
\toprule
Method & F$_{\max}$ & mAUPR & S-min & Brier & ECE & Viol. \\
\midrule
\multicolumn{7}{l}{\textit{DeepGOCNN scores}} \\
Raw          & 0.554 & 0.039 & 16.30 & 0.00234 & 0.000938 & 0.0397 \\
Max closure  & 0.622 & 0.038 & 16.95 & 0.00269 & 0.001758 & 0 \\
Loopy BP     & 0.654 & 0.041 & 15.13 & 0.00202 & 0.001247 & 0 \\
Exact WMC    & 0.655 & 0.041 & 15.13 & \textbf{0.00202} & \textbf{0.001234} & \textbf{0} \\
\addlinespace
\multicolumn{7}{l}{\textit{MLP scores}} \\
Raw          & 0.586 & 0.179 & 14.22 & 0.00216 & 0.001048 & 0.0681 \\
Max closure  & 0.655 & 0.178 & 15.66 & 0.00236 & 0.001614 & 0 \\
Loopy BP     & 0.670 & 0.164 & 14.02 & 0.00195 & 0.001337 & $1.5{\times}10^{-7}$ \\
Exact WMC    & 0.671 & 0.163 & 14.01 & \textbf{0.00195} & \textbf{0.001324} & \textbf{0} \\
\bottomrule
\end{tabular}
\caption{Post-processing DeepGOCNN and MLP scores on complete GO cc
under the CAFA protocol. Higher is better for F$_{\max}$ and mAUPR;
lower is better for all others. Exact WMC achieves zero violations and
the lowest Brier/ECE in both settings.}
\label{tab:downstream}
\end{table}

% ============================================================
\section{Discussion}
\label{sec:discussion}
% ============================================================

\subsection{Why Biomedical Ontologies Are Amenable to Modular Decomposition}

The success of modular separator WMC on GO, UBERON, and ChEBI rests on
a structural property of these ontologies: despite their large size and
high treewidth in the monolithic graph, a small fraction of high-degree
variables (the core) acts as a separator that isolates the remaining
terms into many small, tractable peripheral modules. Table~\ref{tab:structure}
shows that cores of 2--9\% of all terms suffice to produce modules
covering 91--100\% of terms at boundary-last widths of 5--12. This
tree-like periphery with a narrow bottleneck is characteristic of
biomedical ontologies built around hierarchical subsumption: most terms
have few parents, and reconvergence is concentrated at a small number
of high-degree terms.

This structural property is not shared by all knowledge bases. Ontologies
with dense cross-concept constraints (e.g., disjointness or covering
axioms at scale) may produce larger cores and wider boundaries, reducing
the benefit of modular decomposition. Section~\ref{sec:tt-separator}
already handles wide boundaries through TT separators; the limits of
this approach are explored in \suppref{app:nesy}.

\subsection{The Role of Joint Boundary Factors}

A key design choice is the use of \emph{joint} separator factors
$F_j(y_{B_j})$ rather than factored approximations. The reconvergence
analysis in Section~\ref{sec:local-reconvergence} shows why factored
approximations fail: when multiple paths from a common ancestor
converge at a shared boundary variable, the joint distribution over
boundary variables encodes the non-independence induced by the
eliminated module. Table~\ref{tab:direct-tt} confirms empirically that
the product baseline (which uses a factored separator) consistently
produces substantially higher errors than TT-2, even though both
respect the logical support.

\subsection{Limitations}

\paragraph{GO cc degeneracy.}
The current experimental design validates exactness on GO cc
(Table~\ref{tab:exactness}) but the GO cc decomposition reduces to
monolithic JT. We address this by pointing to the GO bp exact-reference
modules in Table~\ref{tab:direct-tt}, which show machine-precision
recovery on genuine modular instances. Future work should produce a
standalone table of exactness results for GO mf and GO bp subgraphs
with depth$\geq 1$ decompositions, as outlined in the experiment
design of Appendix~\ref{app:future-experiments}.

\paragraph{Approximate separators lack a posteriori certificates.}
Proposition~\ref{prop:boundary-error} bounds module marginal error by
boundary posterior error, but does not provide a computable
certificate for $\mathrm{TV}(\widehat{p}(y_{B_j}), p_T(y_{B_j}))$
when exact references are unavailable. Rank self-consistency
(Table~\ref{tab:rank-consistency}) provides indirect evidence that
rank 8 is near saturation; but for safety-critical applications,
certified error bounds would be needed.

\paragraph{Scalability at rank $>$ 8.}
For ChEBI with max $|B_j|=1{,}082$, rank-8 TT separators achieve
median validation NMAE $9.7\times10^{-3}$, which may be insufficient
for tasks requiring high marginal accuracy. Higher ranks are feasible
in principle but increase bond-state cost in reduced-core inference.
Adaptive rank selection and hierarchical TT representations are
directions for future work.

\paragraph{Coverage.}
Table~\ref{tab:operational-scale} reports structural coverage of
91--100\%: between 0--9\% of terms lie in the core rather than in
peripheral modules. Core-variable marginals are obtained from
reduced-core inference and are exact (or as accurate as the TT
separator quality). The remaining terms' marginals are within the
coverage fraction.

\subsection{Future Directions}

\paragraph{Downstream evaluation on GO mf and GO bp.}
The CAFA downstream experiment (Section~\ref{sec:rq5-downstream}) is
on GO cc, where exact JT inference is feasible. Extending this
evaluation to GO mf and GO bp---where TT-mode modular WMC is
the only tractable approach---would directly demonstrate the practical
value of the wide-separator construction.

\paragraph{Adaptive decomposition and rank selection.}
The current core-selection algorithm uses static degree scores.
Score-dependent decomposition (adapting the core to the predictor score
vector at each inference call) could improve module tractability.
Similarly, per-module rank selection based on boundary structure could
reduce TT cost without sacrificing accuracy.

\paragraph{Extensions to richer ontology languages.}
The framework currently handles is-a implications (Horn 2-clauses).
Extensions to OWL DL axioms (covering, disjointness, existential
restrictions) would require handling non-monotone clauses in the module
factor graph, changing the oracle interface and TT support structure.

% ============================================================
\section{Conclusion}
\label{sec:conclusion}
% ============================================================

We introduced modular separator WMC, a framework for scaling
ontology-conditioned weighted model counting to large biomedical
ontologies. The core technical contributions are: (i)~a
boundary-aware core-selection algorithm that decomposes the constraint
graph into a small core and tractable peripheral modules; (ii)~exact
and nonneg\-ative TT separator constructions that summarize each
module's effect on the rest of the ontology through a joint boundary
factor; and (iii)~backward calibration that recovers all module
marginals from saved junction-tree messages.

Two propositions ground the framework theoretically. Exact separators
provably recover the monolithic WMC posterior, and module marginal
error is bounded by boundary posterior error under total-variation
distance. These guarantees hold regardless of how the boundary
posterior is computed, separating the approximation question from the
recovery correctness question.

Experiments on five complete biomedical ontologies show that
decomposition enables all-term inference beyond the 300\,GB memory
limit of monolithic junction-tree inference, that state-of-the-art WMC
solvers exhaust memory on the same inputs, and that rank-controlled TT
separators systematically improve accuracy as rank increases. On the
CAFA benchmark, exact posterior correction eliminates true-path
violations and improves calibration over raw predictor scores and loopy
belief propagation. The results show that ontology structure---sparse
DAG with a small high-degree core and many tractable periphery
modules---is precisely the regime where modular decomposition and
joint separator factors are most effective.

% ============================================================
\bibliography{references}
% ============================================================

% ============================================================
\appendix
% ============================================================

\section{Future Experiment Designs}
\label{app:future-experiments}

To fully address reviewer c36U's concern about GO cc degeneracy, the
following experiment provides standalone exactness validation on
non-trivial decompositions. For GO mf and GO bp, sample subgraphs
(depth-6 neighborhoods of random roots) retaining instances where
Algorithm~\ref{alg:core-selection} produces depth$\geq 1$ and at least
one module with $\widehat{w}_j^\partial\leq 12$. Collect $\geq 100$
such instances per ontology, generate 4 score vectors each, and compare
modular WMC (exact mode) against brute-force and monolithic JT
references. Expected results: MAE$\approx 10^{-15}$ across all
instances, matching Proposition~\ref{prop:exact-modular}.

For GO mf CAFA downstream evaluation, map CAFA3 annotations to the GO
mf namespace and compare TT-mode modular WMC ($r=8$) against loopy BP
and max-closure using the same CAFA metrics. This directly demonstrates
the practical value of the wide-separator construction on a namespace
where the decomposition is genuinely required.

\section{Reproducibility Details}
\label{app:reproducibility}

The operational selector uses Algorithm~\ref{alg:core-selection} with
$w_{\max}=14$, $\beta=0.01$, and $\rho=0.30$. Direct TT fitting uses
$r\in\{1,2,4,8\}$, 5,000 unique training queries, six
alternating-least-squares sweeps, nonneg\-ative least-squares core
updates, and an independent validation set. Scores in TT mode are
clipped to $[10^{-6},1-10^{-6}]$. Loopy BP uses synchronous
sum-product with damping 0.5, tolerance $10^{-9}$, and at most 20,000
iterations. Exact-reference candidates retain 108 subgraphs per
ontology where brute-force and JT references agree under fixed random
seeds.

\paragraph{Compute resources.}
CPU-only SLURM runs. Multi-ontology screening uses 16 CPU cores,
128\,GB allocated memory. The full ChEBI recursive run uses 1 CPU core,
400\,GB allocated memory. Monolithic JT rows use a common 300\,GB cap.

\paragraph{Data availability.}
Source code, ontology preprocessing, exact and approximate inference
backends, Lean developments, environment files, seeds, raw results, and
table-generation scripts are included in the submission artifact.

\section{Rank Self-Consistency (Extended)}
\label{app:rank-consistency}

Table~\ref{tab:rank-consistency} reports median cross-rank MAE for ranks
2, 4, and 8. For GO mf and GO bp, rank-16 runs confirm the convergence
trend: rank-8 to rank-16 MAE is $3.7\times10^{-3}$ (GO mf) and
$5.2\times10^{-3}$ (GO bp), consistent with the rank-4 to rank-8
reduction. The saturation pattern holds across all four ontologies.

\section{Additional Theory and Diagnostics}
\label{app:additional-theory}

\subsection{General Boolean Queries}

For any Boolean query $\psi$,
\[
  p_T(\psi) = \frac{\mathrm{WMC}(T\wedge\psi,p)}{\mathrm{WMC}(T,p)}.
\]

\subsection{Star Exactness}

\begin{theorem}[Star exactness]
\label{thm:star}
For a parent $\pi$ with independent children $c_1,\ldots,c_k$ under
constraints $y_{c_j}\Rightarrow y_\pi$, the exact WMC marginal is
\[
  \mu_\pi = \frac{p_\pi}{p_\pi + (1-p_\pi)\prod_{j=1}^k(1-p_{c_j})}.
\]
\end{theorem}

\begin{proof}
Condition on $y_\pi$. When $y_\pi=1$, all child assignments are
allowed, contributing $p_\pi\prod_j[p_{c_j}+(1-p_{c_j})] = p_\pi$.
When $y_\pi=0$, every child must be false, contributing
$(1-p_\pi)\prod_j(1-p_{c_j})$. Taking the ratio gives the stated formula.
\end{proof}

\subsection{Diamond Calculation}
\label{app:diamond-calculation}

For $T_\diamond$ with all four Bernoulli parameters equal to $1/2$,
six assignments satisfy the implications; five have $y_a=1$, giving
$p_{T_\diamond}(y_a=1) = 5/6$. The path-wise calculation gives
$\widehat\mu_a = 9/10$ (Section~\ref{sec:local-reconvergence}).

\subsection{Dense-Mode Cost}
\label{app:dense-cost}

Exact dense-mode elimination costs
\[
  O\!\left(
    \sum_{j=1}^s\!\left[n_j 2^{\widehat{w}_j^\partial+1}+2^{|B_j|}\right]
    + (|C|+s)2^{\widehat{w}_C+1}
  \right).
\]

\section{Ontology Instantiation Details}
\label{app:multi-ontology}

GO experiments: release 2021-11-16.
UBERON: \texttt{uberon-basic.owl} release 2026-06-19.
ChEBI: OBO data-version 252, dated 2026-05-01.
Source paths, URLs, and SHA256 checksums are recorded in the artifact
manifest.

\section{Structure Decomposition Table}
\label{app:structure}

\begin{table*}[t]
\centering
\small
\begin{tabular}{lrrrrrrrr}
\toprule
Ontology & Terms & \%mp & Max deg.\ & $\widehat{w}(G_T)$
& Core & Modules & $\max_j\widehat{w}(A_j)$ & Struct.\ cov. \\
\midrule
GO cc  & 4,041   & 13.0 & 4  & 12  & 0.0\% & 1       & 12 & 100.0\% \\
GO mf  & 10,131  & 20.0 & 7  & 37  & 5.0\% & 6,519   & 5  & 95.0\%  \\
GO bp  & 25,950  & 53.0 & 9  & 269 & 9.0\% & 7,249   & 12 & 91.0\%  \\
UBERON & 14,977  & 24.0 & 7  & 46  & 5.0\% & 8,427   & 11 & 95.0\%  \\
ChEBI  & 205,592 & 24.7 & 14 & 738 & 2.0\% & 181,679 & 12 & 98.0\%  \\
\bottomrule
\end{tabular}
\caption{Unrestricted modular pre-screen. Structural coverage is
$\sum_j|M_j|/|Y|$.}
\label{tab:structure}
\end{table*}

\section{Formalisation: Lean~4 Development}
\label{app:lean}

The Lean~4 development (\texttt{Wmc.lean}, Lean 4.31) is a
single Lean-core file with no dependence on \texttt{mathlib}. It
defines worlds as Boolean lists, the true-path constraint as a
conjunction of edge clauses $(\neg y_c\vee y_\pi)$, and closes every
claim by \texttt{decide} using kernel-checkable rational arithmetic.
Verified lemmas include \texttt{diamond\_exact\_5\_6} ($5/6$),
\texttt{diamond\_softor\_9\_10} ($9/10$), and
\texttt{diamond\_softor\_ne\_exact} (inequality). The parametric star
theorem is proved in a separate \texttt{mathlib} development with
axiom footprint $\{\texttt{propext},\texttt{Classical.choice},
\texttt{Quot.sound}\}$.

\end{document}
